\documentclass[12pt,a4paper]{article}

% ============================================================
% PACKAGES
% ============================================================
\usepackage[utf8]{inputenc}
\usepackage[T1]{fontenc}
\usepackage{lmodern}
\usepackage[english]{babel}
\usepackage{microtype}
\usepackage{amsmath}
\usepackage{geometry}
\usepackage{setspace}
\usepackage[hidelinks]{hyperref}
\usepackage{indentfirst}

\geometry{top=2.5cm, bottom=2.5cm, left=2.5cm, right=2.5cm}
\onehalfspacing
\setlength{\parindent}{1.25cm}
\setlength{\parskip}{0pt}
\setcounter{section}{1}

\begin{document}

% ============================================================
% 2. MATERIALS AND METHODS
% ============================================================
\section{Materials and methods}
\label{sec:methods}

\subsection{Animals and ethical approval}
\label{sec:animals}
Adult wild-type AB zebrafish (\textit{Danio rerio}) were obtained from a commercial supplier (Angelfish, Londrina, Paran\'a, Brazil). Fish were of mixed sex, had no previous experimental use, and had a mean total length of $3.4 \pm 0.4$~cm and a mean body mass of $0.3 \pm 0.1$~g. A total of 410 fish were analysed across the three testing contexts after the exclusions described in Section~\ref{sec:data-preparation}. All procedures were approved by the Animal Use Ethics Committee of the State University of Londrina (CEUA/UEL; authorisation 040.2020).

\subsection{Husbandry and acclimation}
\label{sec:husbandry}
Fish were acclimated for at least 10 days before testing. They were housed in approximately 70~L glass aquaria containing dechlorinated tap water at a maximum density of 1~g of fish per litre. Water temperature was maintained at $26 \pm 1\,^{\circ}\mathrm{C}$ under a 14:10~h light--dark photoperiod. Water conditions were held at pH $7.0 \pm 0.5$, conductivity $110 \pm 10~\mu\mathrm{S\,cm^{-1}}$, nitrite $\leq 0.25$~mg/L, and non-ionised ammonia $\leq 5~\mu$g/L. Aquaria were fitted with sponge filters and aerators. Debris was removed by siphoning when necessary, 50--75\% of the water was renewed weekly, and fish received a commercial diet twice daily.

\subsection{Study design and analytical hierarchy}
\label{sec:study-questions}
The experiment crossed three nominal ethanol concentrations---0\% (control), 0.5\%, and 1\%---with three exposure durations of 1, 24, and 96~h. Each fish was allocated to one concentration, one exposure duration, and one testing context, and each concentration-by-duration-by-context cell contained 12--16 fish.

The study was organised around three questions. The primary question was methodological: whether measurements obtained when an assay occupied a position within the sequential battery differed from measurements obtained in the corresponding comparison context. The exploratory question was integrative: how behavioural, endurance, physiological, and condition measurements obtained from the same fish related to one another. The secondary question was pharmacological: how ethanol concentration and exposure duration affected the modelled endpoints within the sequential battery.

This hierarchy applies to the study questions, not to a ranking of endpoints. No endpoint was designated a primary endpoint. Endpoints were entered in a registry according to their domain, data type, and analytical role, and the four analyses---battery-context comparisons, ethanol endpoint models, correlations, and principal component analysis---were kept separate because they answer different questions.

One term recurs throughout the analysis and is defined here. A \textit{comparison cell} in the battery-context analysis is one endpoint measured at one ethanol concentration and one exposure duration, compared between a specified pair of testing contexts. Thirteen behavioural and performance endpoints were eligible for battery-context comparison: five from the light-dark test, four from the novel tank test, and four from the swimming-resistance assay. With nine concentration-by-duration combinations per endpoint, this yields 117 comparison cells. The 117 comparisons therefore represent endpoint-by-concentration-by-duration cells, not 117 distinct endpoints.

\subsection{Ethanol exposure}
\label{sec:ethanol-exposure}
Exposure media were prepared with absolute ethanol (99.9\%). Before ethanol was added, an equivalent volume of culture water was removed so that the final volume remained constant; control media contained culture water without ethanol. Concentrations are reported as nominal because analytical confirmation of the achieved concentrations was not available.

Fish assigned to the 1~h exposure were placed individually in 250~mL vessels containing 200~mL of the assigned solution for 30~min before testing. In the sequential battery this vessel time was followed by the light-dark test, comprising 5~min of acclimation and 10~min of observation, bringing cumulative ethanol exposure to 45~min, and then by the 15~min novel tank test, bringing cumulative exposure to 60~min. Ethanol at the assigned concentration was present in the holding vessel and in the light-dark and novel tank media. The swimming-resistance assay was conducted in culture water without ethanol, because its large recirculating system could not maintain a stable ethanol concentration.

Fish assigned to the 24 and 96~h exposures were exposed collectively in 30~L aquaria to which they had been acclimated for at least 5 days. Filters were removed and aeration was reduced after ethanol was added. At the end of the assigned exposure period, fish were transferred individually to vessels containing 200~mL of the same exposure medium before behavioural testing. As in the 1~h groups, ethanol exposure continued through the light-dark and novel tank tests, while the swimming-resistance assay was conducted in culture water.

Holding times in the comparison contexts were set so that each assay began at the same cumulative exposure time at which that assay occurred in the sequential battery. Fish in the reversed behavioural context were held for 45~min before the novel tank test and then completed the light-dark test; fish in the isolated endurance context were held for 60~min before the swimming-resistance assay. These timing arrangements form part of the measurement context and are not interpreted as order effects in isolation.

\subsection{Testing contexts and battery architecture}
\label{sec:design}
Three testing contexts were used (Figure~1; Table~1). Because different fish were tested in each context, every context comparison is a between-group comparison. It reflects assay position together with elapsed experimental time, preceding handling and activity, and any cohort-related differences, and does not isolate a within-animal carry-over effect.

\subsubsection{Sequential battery (BT1)}
\label{sec:bt1}
BT1 was the complete sequential battery and included 136 fish. Each fish completed the light-dark test, then the novel tank test, and finally the swimming-resistance assay. The order was fixed so that the approach--avoidance assay preceded accumulated handling and physical exertion, and the most physically demanding assay came last. The battery was run as an overlapping workflow: as one fish advanced to the next station, the preceding apparatus became available for another fish, so that all three stations could operate concurrently while every individual still contributed measurements from all three domains.

\subsubsection{Reversed behavioural context (BT2)}
\label{sec:bt2}
BT2 included 139 fish and reversed the order of the first two assays, with the novel tank test followed by the light-dark test. The swimming-resistance assay was not performed. BT1 and BT2 therefore provide position-specific comparisons rather than a comparison between one sequential and one isolated group. For the light-dark test, BT1 provides the first-position measurement and BT2 the measurement taken after the novel tank test. For the novel tank test, BT2 provides the first-position measurement and BT1 the measurement taken after the light-dark test. The first-position assay in either context had no preceding behavioural assay, but remained part of its assigned experimental context.

\subsubsection{Isolated endurance context (BT3)}
\label{sec:bt3}
BT3 included 135 fish and consisted of the swimming-resistance assay alone. It provides the isolated comparator for the endurance measurements that follow the light-dark and novel tank tests in BT1.

\subsection{Behavioural and performance assays}
\label{sec:assays}

\subsubsection{Light-dark test}
\label{sec:ldt}
The light-dark test was conducted in a glass aquarium measuring $18 \times 49 \times 9$~cm (height $\times$ length $\times$ width) containing 6.8~L of solution, divided into equally sized light and dark compartments by white and black adhesive film applied to the outer walls. Each fish was placed in a central compartment for 5~min of acclimation. A remotely operated partition was then raised and behaviour was recorded for 10~min with an overhead camera.

Five endpoints were obtained:
\begin{itemize}
    \item first choice: the first compartment entered after release, coded 0 for dark and 1 for bright;
    \item latency: the time from release to the first compartment choice;
    \item time in the bright zone, in seconds;
    \item the number of transitions between compartments; and
    \item mean time per crossing (MTPC), calculated as time in the bright zone divided by the number of transitions.
\end{itemize}
MTPC therefore expresses mean bright-zone time per transition; it is not the time taken to cross between compartments, and it is mathematically related to both bright-zone time and transition count. Fish that did not make a valid first choice were coded as missing for the binary first-choice endpoint.

\subsubsection{Novel tank test}
\label{sec:ntt}
The novel tank test used a trapezoidal apparatus measuring $14 \times 24.3 \times 4.5$~cm containing 1.5~L of solution, with matte acrylic false walls forming the trapezoidal profile and reducing reflections. Behaviour was recorded for 15~min with a front-mounted camera, and trajectories were analysed with SACAM video-analysis software (GeoVision). Endpoint calculations used a predefined 480-s observation window following acclimation, with the tank divided into upper and lower strata.

Four endpoints were calculated:
\begin{itemize}
    \item total distance travelled, in centimetres;
    \item total movement time, the observation window minus periods without detected movement;
    \item mean velocity, total distance divided by movement time; and
    \item upper-stratum preference, the percentage of movement time spent in the upper half of the tank.
\end{itemize}
The experimental records establish the 15-min recording period, the acclimation period, and the 480-s analytical window, but do not identify the placement of the remaining approximately 2~min within the recording; no more precise placement was assumed.

\subsubsection{Swimming-resistance assay}
\label{sec:resistance}
Swimming resistance was measured in a closed recirculating apparatus that exploits positive rheotaxis. A pump generated a manually adjustable current that passed through an internal mesh into a transparent acrylic tube of 18~mm internal diameter, connected to an escape chamber fitted with a heater and thermometer. Fish entered through an opening at the top of the apparatus, and a lid at the base prevented their return once testing began. Approximately 75\% of the medial portion of the tube was covered with aluminium foil to reduce external visual cues. Flow began at 1~L/min and increased by 1~L/min every minute to a maximum of 12~L/min. A fish that failed to respond to the initial flow and was carried into the escape chamber was reinserted into the tube, to a maximum of five attempts.

The resistance index was calculated as the sum of the flow rates sustained for a complete minute, plus the final flow rate multiplied by the fraction of the final minute sustained:
\begin{equation}
\label{eq:resistance-index}
\begin{split}
RI = \sum_{i=1}^{k-1} Q_i \;+\; Q_k\left(\frac{t_k}{60}\right),
\end{split}
\end{equation}
where $Q_i$ is the flow rate of each completed one-minute step, $Q_k$ is the final flow rate reached, and $t_k$ is the number of seconds sustained at that final flow rate. For example, a fish that completed the steps from 1 to 9~L/min and then remained for 30~s at 10~L/min received an index of:
\begin{equation}
\label{eq:resistance-example}
\begin{split}
1+2+3+4+5+6+7+8+9 \;+\; 10 \times \frac{30}{60} \;=\; 50.
\end{split}
\end{equation}

Lower values indicate a reduced capacity to swim against the current. Four endurance endpoints were retained: the resistance index; the final flow step reached, an ordinal variable with levels 1--12; the number of reinsertion attempts; and the number of seconds sustained at the final flow step. Because the resistance index incorporates both flow progression and the final step reached, it and the final flow step are related measurements and were not treated as independent confirmation of performance.

\subsection{Post-test physiology and morphometrics}
\label{sec:physiology}

\subsubsection{Blood glucose}
\label{sec:glucose}
After testing, fish were anaesthetised by immersion in ice-cold water until opercular movement ceased. A caudal section was made with surgical-steel scissors and blood was collected immediately. Glucose was measured with an Abbott FreeStyle Optium Neo portable glucometer and recorded in mg/dL.

Because the fish were small relative to the collection procedure, the blood volume obtained was sometimes insufficient to complete the measurement. Such observations were recorded as missing rather than estimated. Glucose was unavailable for 25 of 136 BT1 fish, 102 of 139 BT2 fish, and 62 of 135 BT3 fish. Because this missingness was substantial and differed among contexts, blood glucose was not included in the battery-context comparisons. It was retained in the ethanol endpoint model within BT1 and in the pairwise-complete BT1 correlation analysis.

\subsubsection{Morphometrics and condition factor}
\label{sec:condition}
Body mass was measured with a digital balance to two decimal places. Total and standard lengths were measured with a caliper, standard length extending from the rostral tip to the base of the caudal-fin rays. Sex was determined from hooks or rugosities on the pectoral-fin rays, which occur in males. Fulton's condition factor was calculated from the audited operational variables as:
\begin{equation}
\label{eq:condition-factor}
\begin{split}
K = 100 \left( \frac{\mathrm{body\ mass\ (g)}}{\mathrm{standard\ length\ (cm)}^{3}} \right).
\end{split}
\end{equation}

Condition factor was treated as a descriptive control measurement. It was included in the exploratory correlation analysis and the principal component analysis, but was not among the ethanol-modelled endpoints and was not used in the battery-context comparisons.

\subsection{Endpoint registry, data preparation, and exclusions}
\label{sec:data-preparation}
The endpoint registry contained 15 measurements: five from the light-dark test, four from the novel tank test, four from the swimming-resistance assay, blood glucose, and condition factor (Table~2). Different analytical sets were used because each analysis addresses a different question:
\begin{itemize}
    \item the battery-context analysis used the 13 behavioural and endurance endpoints (five light-dark, four novel tank, four resistance);
    \item the ethanol endpoint models used 14 endpoints, namely those 13 plus blood glucose;
    \item the correlation analysis used all 15 registered measurements, including blood glucose and condition factor; and
    \item the principal component analysis used 14 endpoints, namely the 13 behavioural and endurance endpoints plus condition factor, with blood glucose excluded to preserve the complete-case sample.
\end{itemize}

The initial dataset contained 414 fish. Four were excluded globally: two from BT1, comprising one light-dark protocol violation and one death during testing; one from BT2 with entirely missing novel tank data; and one from BT3 with entirely missing endurance data. The analysed samples were therefore BT1 $n=136$, BT2 $n=139$, and BT3 $n=135$, giving 410 fish in total.

First choice was missing for one BT1 fish and one BT2 fish under the no-choice or ambiguous-choice rule. Endpoint-specific missing values were retained as missing and handled according to the requirements of each analysis. No values were imputed. Derived variables were calculated only according to the formulas specified above.

\subsection{Statistical analysis}
\label{sec:statistics}
Descriptive summaries included sample size, mean, standard deviation, median, interquartile range, minimum, and maximum where appropriate. Inference used a two-sided threshold of $\alpha=0.05$, and Benjamini--Hochberg adjustment controlled the false discovery rate within the relevant analytical families; adjusted $p$-values are denoted $p_{\mathrm{BH}}$. Battery-context comparisons, ethanol endpoint models, correlations, and principal component analysis were analysed separately, and the correlation and component analyses were descriptive and were not used as inferential evidence of ethanol or measurement-context effects.

\subsubsection{Primary battery-context comparisons}
\label{sec:context-statistics}
Battery-context comparisons were made within matched concentration-by-duration cells. BT1 was compared with BT2 for the five light-dark and four novel tank endpoints, giving
\begin{equation}
\label{eq:bt1-bt2}
\begin{split}
9\ \mathrm{endpoints} \times 9\ \mathrm{cells} = 81\ \mathrm{comparisons},
\end{split}
\end{equation}
and BT1 was compared with BT3 for the four endurance endpoints, giving
\begin{equation}
\label{eq:bt1-bt3}
\begin{split}
4\ \mathrm{endpoints} \times 9\ \mathrm{cells} = 36\ \mathrm{comparisons}.
\end{split}
\end{equation}
The primary analysis therefore comprised 117 endpoint-by-concentration-by-duration comparisons.

Non-binary endpoints were compared with Mann--Whitney rank-sum tests, and first choice was compared with Fisher's exact test. All 117 raw $p$-values were adjusted together with the Benjamini--Hochberg procedure, and a context difference was classified as detected only when the adjusted $p$-value fell below 0.05.

\subsubsection{Effect sizes and uncertainty}
\label{sec:effect-sizes}
Effect sizes accompanied every comparison, because the size of a difference and the precision with which it is estimated are more informative than a threshold decision. For non-binary endpoints we report Hedges' $g$, rank-biserial correlation, and Cliff's delta. Hedges' $g$ expresses the difference between contexts relative to the pooled within-cell variability and includes a correction that makes it suitable for the cell sizes used here; we do not attach fixed verbal categories to particular values, since what counts as a meaningful difference depends on the endpoint. For first choice, risk differences and odds ratios are reported.

Raw-scale mean, median, and risk differences were accompanied by percentile bootstrap confidence intervals based on 2,000 resamples, which express the range of values compatible with the data rather than a decision rule. All estimates are oriented as the comparison context minus BT1, so a positive estimate indicates higher measurements in BT2 or BT3 than in BT1 and a negative estimate indicates lower measurements in the comparison context.

\subsubsection{Equivalence boundary}
\label{sec:equivalence}
The battery-context analysis tested whether the data provided evidence of a context difference. It did not test whether measurements were sufficiently similar to be treated as equivalent. Formal equivalence requires an endpoint-specific smallest effect size of interest to be justified before the results are examined, with equivalence supported only when the estimated effect and its uncertainty fall within those prespecified bounds, commonly assessed with two one-sided tests. No biologically or operationally justified equivalence margins were specified prospectively in this study, and equivalence testing was therefore not conducted.

A comparison without a detected adjusted difference is accordingly described as showing no Benjamini--Hochberg-adjusted context difference, and not as demonstrating that the measurements are equivalent or interchangeable. Prospective equivalence studies should define endpoint-specific margins, justify them independently of the observed effects, and power the experiment to evaluate those margins, following the framework described by Lakens et al.\ (2018).

\subsubsection{Secondary ethanol endpoint models}
\label{sec:ethanol-statistics}
Ethanol responses were analysed within BT1 only, each endpoint being modelled as a function of ethanol concentration, exposure duration, and their interaction. The model family was matched to the measurement scale of each endpoint.

Time in the bright zone, MTPC, total distance, mean velocity, movement time, upper-stratum preference, resistance index, time at the final flow step, and blood glucose were analysed with aligned-rank-transform analysis of variance, which provides rank-based inference rather than effect estimates in the original units. Zone-transition count was analysed with a negative-binomial generalised linear model and the number of attempts with a Poisson generalised linear model. First choice was analysed with Firth's penalised logistic regression because of sparse or separated cells, the no-choice BT1 observation being excluded from this model; coefficients were exponentiated to obtain odds ratios and confidence intervals.

Latency contained 98 zero values among the 136 BT1 fish and was analysed with a two-part hurdle model. The first component used logistic regression for the probability of a latency greater than zero, and the second used a log-normal linear model for $\log(\mathrm{latency})$ among the 38 fish with a positive latency, with treatment, exposure, and interaction terms in both components. The hurdle model is the governing latency analysis and supersedes the aligned-rank result generated before zero inflation was addressed. The ordinal final flow step was analysed with proportional-odds logistic regression over ordered levels 1--12, with treatment, exposure, and interaction terms evaluated by type-II likelihood-ratio tests.

Benjamini--Hochberg adjustment was applied within the corresponding model or coefficient families. Where a relevant aligned-rank term was detected, estimated marginal contrasts were calculated with the same adjustment and summarised as compact-letter displays. These displays represent comparisons on the aligned-rank scale and do not directly order measurements in their original units, so raw medians are reported alongside them. Ethanol-effect inference rests exclusively on these endpoint-level models, and not on the component analysis or on visual separation among groups.

\subsubsection{Exploratory correlation analysis}
\label{sec:correlations}
Relationships among the 15 registered BT1 measurements were described with Spearman rank correlations using pairwise-complete observations. Pairwise sample sizes ranged from 110 to 136 because blood-glucose availability varied. Correlations were calculated across the pooled BT1 fish rather than within concentration-by-duration cells, and were used to identify convergence, separation, and redundancy among measurements. Nominal $p$-values annotate the display descriptively and were not used to infer treatment effects, and no latent constructs were inferred from the pattern of association.

Correlations involving derived or operationally linked variables were interpreted with that dependence in mind. MTPC contains bright-zone time in its numerator, mean velocity is calculated from distance and movement time, and the resistance index incorporates flow progression and the final flow step; associations among such measurements were not treated as independent confirmation of the same effect.

\subsubsection{Exploratory principal component analysis}
\label{sec:pca}
A single principal component analysis was used to visualise covariance among 14 standardised BT1 endpoints: time in the bright zone, MTPC, first choice, latency, zone transitions, total distance, mean velocity, movement time, upper-stratum preference, resistance index, final flow step, attempts, time at the final flow step, and condition factor.

Blood glucose was excluded because requiring complete glucose data would have reduced the analysis to 110 fish. Excluding it retained 135 of the 136 BT1 fish, the single omission being the fish with the missing first-choice value. No values were imputed. Blood glucose remained in its endpoint model and in the pairwise-complete 15-variable correlation analysis, so the difference between the 15-variable correlation analysis and the 14-endpoint component analysis is deliberate.

Variables were centred and scaled to unit variance, and the analysis was computed by singular-value decomposition equivalent to \texttt{prcomp} with centring and scaling. The two-dimensional biplot displays PC1, explaining 23.3\% of variance, and PC2, explaining 17.2\%, for a cumulative 40.6\%, with all 14 loading vectors retained. PC3 explained 10.3\% and is reported in the supplementary scree plot and the optional interactive three-dimensional display, without being assigned an additional interpretation.

Because the analysis represents standardised continuous, binary, count, bounded, and ordinal measurements, and because several variables are mathematically or operationally related, it was used strictly as exploratory visualisation. Components are described from their loading relationships and are not treated as measured latent traits or given fixed conceptual names. Treatment and exposure symbols appear for orientation only, and no ethanol or context inference is drawn from the arrangement of scores. The interactive three-dimensional display uses the same archived canonical scores and loadings as the two-dimensional biplot; it does not refit the analysis, alter loadings, or add an inferential result.

\subsection{Software and reproducibility}
\label{sec:software}
Data preparation and auditing used Python and R, and inferential analyses were conducted in R version 4.5.3. Aligned-rank-transform models used ARTool version 0.11.2, estimated marginal means used emmeans version 2.0.2, likelihood-ratio and omnibus tests used car version 3.1-5, proportional-odds and negative-binomial models used MASS, Firth logistic regression used logistf, and bootstrap intervals used boot version 1.3-32. The canonical principal component analysis was computed from the centred and scaled matrix by a \texttt{prcomp}-equivalent singular-value decomposition, and the optional three-dimensional display was generated only from the archived canonical scores and loadings. Statistical tables and data visualisations were produced from the audited data or saved model outputs through reproducible computational procedures.

\subsection{Use of generative artificial intelligence}
\label{sec:ai-disclosure}
During the preparation of this work the authors used OpenAI GPT Codex Sol 5.6, High setting, and Anthropic Claude Code Opus 5.0, High setting. These systems assisted with manuscript organisation, drafting and language revision, inspection of analysis code, consistency checks among scripts, saved outputs, tables, figures, and manuscript text, literature organisation, and documentation of the analytical workflow.

All AI-assisted material was checked by the human author against the source data, the saved analyses, the analysis code, the original thesis, and the source literature. No result, reference, or endpoint definition was accepted solely from model-generated text. The systems did not generate raw observations, determine the experimental design, or perform or approve the analyses. Final scientific, methodological, statistical, interpretive, and editorial decisions remained with the human author, who accepts responsibility for the accuracy and integrity of the manuscript. The tools are not listed as authors because they cannot take responsibility for the work.

\end{document}
